\subsection{Introducción}
En este documento se definirán las mecánicas y reglas del juego \textit{Magnate}.
Se explicarán las diferentes fases del juego, las acciones que pueden realizar 
los jugadores y las condiciones de victoria. 

Se definirá desde el diseño del tablero hasta los distintos eventos especiales
que pueden ocurrir, considerando todos los casos posibles.

Las reglas están basadas en las reglas originales del Monopoly \cite{hasbro_monopoly_es} 
tomando inspiración adicional de la saga de videojuegos Mario Party.
Magnate busca ser un juego dinámico con alto grado de aleatoriedad, pero
sin perder la estrategia del Monopoly original.

\subsection{Objetivo del juego}
El objetivo del juego es ser el jugador con mayor valor neto al final de la
partida, limitada por un número de rondas. El valor neto de un jugador se calcula
mediante la suma de su dinero en efectivo más el valor de venta
de todas sus propiedades, sin contar las que tenga hipotecadas y
el valor de venta de sus construcciones.



\subsection{Equipamiento}
A continuación está la lista de elementos para una partida.
\begin{itemize}
	\item Tablero
	\item 2 Dados normales
	\item 1 Dado bus 
	\item Cartas de las propiedades 
	\item Casas
	\item Hoteles 
	\item Cartas "Sales de la cárcel"
	\item Cartas de evento fantasía
\end{itemize}

\subsubsection{Los dados: definición y funcionamiento}
En el juego hay 3 dados en juego, 2 normales de 6 caras numerados del 1 al 
6 y uno especial llamado \textit{Dado bus}. El Dado bus está compuesto de 6
caras, 3 de ellas con el dibujo de un bus y las otras 3 numeradas del 1 al 3.

En una tirada de dados se tiran siempre los 3 dados a la vez.
En caso de salir 3 números se tendrá en cuenta para el movimiento
la suma total de los 3. En caso de salir un bus el jugador podrá elegir 3
valores: el valor de uno de los dados normales, el valor del otro de los dados
normales o la suma de los 2 dados normales.

Para las mecánicas de los dobles se tendrán en cuenta únicamente los dados 
normales. Al sacar dobles, el jugador gana una tirada extra. Si llega a 
sacar dobles una tercera vez consecutiva en el mismo turno, va directamente 
a la cárcel. 

Si en la tirada sale un triple (en los casos en los que el Dado bus saque número),
el jugador puede viajar a la casilla del tablero que quiera, aplicándose el efecto
de caer en esa casilla tal y como hubiera sido si la tirada hubiera sido 
convencional. Después de sacar un triple no se vuelve a tirar (como si de unos 
dobles se tratara). Los triples no cuentan en la racha de dobles, por lo que no 
te pueden mandar a la cárcel.

\subsubsection{Tablero}
El tablero está compuesto de 2 pistas cuadradas concéntricas:
el anillo exterior tiene 36 casillas y 4 esquinas y el anillo 
interior tiene 12 casillas y 4 esquinas. La particularidad
de las casillas para transición y cómo afectan al tablero 
se comentará con las casillas de puentes.


\subsection{Casillas}
En esta sección se definirán todos los distintos tipos de casilla 
que hay en el tablero y sus efectos y mecánicas.

\subsubsection{Propiedades}
Las propiedades son casillas que se pueden comprar, vender e hipotecar,
y por las que se cobra alquiler a jugadores rivales. Si un jugador 
cae en una de estas casillas pueden pasar varias cosas en función del 
estado de compra de la propiedad:
\begin{itemize}
	\item La casilla no tiene dueño. En este caso el jugador tiene derecho 
	a comprarla por el precio de compra indicado. En caso de no querer comprar o no 
	tener suficiente dinero, la propiedad pasa a ser subastada.
	\item La casilla tiene dueño. El jugador debe pagar el alquiler calculado
	en función de las construcciones o de si tiene el grupo completo.
	\item La casilla está hipotecada. No ocurre nada.
	\item La casilla es del propio jugador. No ocurre nada.
\end{itemize}

\subsubsection{Grupos de propiedades}
Las propiedades están divididas en grupos de color. Un jugador que posea
todas las propiedades de un mismo grupo, podrá cobrar el doble de alquiler
por las propiedades sin construcciones
y obtendrá derecho a construir en ellas.


\subsubsection{Puentes}
Los puentes son un tipo especial de propiedad. En el juego tienen
una función doble: poder ganar dinero por alquiler (como el resto de 
propiedades), y habilitar la transición entre el cuadrado interior y el 
cuadrado exterior.

El alquiler a pagar en los puentes viene marcado por el número de 
puentes que el jugador posea.

La mecánica de transición viene marcada por el valor usado 
en la tirada de dados (las casillas que el jugador elige mover). Si la 
tirada es par se continúa en el mismo cuadrado, si es impar se cambia de 
cuadrado (o viceversa, lo que indique la casilla).

En caso de caer justo en la casilla del puente, la decisión par o impar 
la tomará la tirada de dados de la siguiente jugada y no la tirada de la jugada 
de caer en el puente.

Las casillas puente, a pesar de ocupar el espacio de 2 casillas convencionales
únicamente contarán como una en los movimientos tanto en los casos de 
mantenerse en el cuadrado como en el caso de transicionar al otro
cuadrado.

\subsubsection{Servidores}
Los servidores son otro tipo especial de propiedad. El valor de su 
alquiler viene determinado por la cantidad de servidores que 
el jugador posea.

\subsubsection{Casillas fantasía}
Al caer en una casilla fantasía, aparecerán 2 cartas, una desvelada y 
una oculta. La carta desvelada tendrá un evento fantasía aleatorio, que se podrá 
aplicar por el precio indicado en la carta. El evento fantasía de la 
carta oculta es gratuito.
En caso de no querer comprar el evento desvelado, o no tener dinero 
para comprarlo, habrá que escoger la carta oculta obligatoriamente.
En caso de comprar y aplicar el evento fantasía de pago, no se aplicará 
ni desvelará el evento oculto.

Los eventos fantasía tendrán acciones positivas (como ganar dinero),
negativas (como perder propiedades) o neutras (como mover a otro 
jugador a cualquier casilla del mapa).

\subsubsection{Vas a la cárcel}
Si caes en esta casilla vas directamente a la cárcel y acaba tu turno.
En el caso de que el jugador deba dinero al banco se le permitirá realizar
las gestiones necesarias para saldar su deuda.

\subsubsection{Cárcel}
Un jugador en la cárcel tiene derecho al principio de su turno a 
tirar los dados. Si saca dobles puede salir de la cárcel usando 
el valor de la tirada y jugar su turno con normalidad (pero sin repetir turno por ser dobles).
Si no saca dobles, puede pagar una fianza de valor fijo. 
Si al tercer turno en la cárcel no ha 
salido, será obligado a pagar la fianza.

Los jugadores en la cárcel tienen derecho a cobrar alquileres y a realizar
acciones de gestión en su turno, pero no pueden moverse por el tablero ni
participar en las subastas.

Los turnos en los que un jugador acabe en la cárcel acaban cuando se 
entre en la cárcel sin tener oportunidad de hacer gestiones (excepto
el caso en el que tenga que saldar deudas con el banco).

\subsubsection{Salida}
En esta casilla comienza la partida y cada vez que un jugador pasa por
ella (al caer justo sobre ella o la primera tirada en la que la dejas atrás
de la vuelta al cuadrado correspondiente) 
recibe una cantidad fija de dinero. 

Si por algún motivo un jugador 
decide mantenerse en el cuadrado interior y no pasar por salida, no recibirá
dinero.

\subsubsection{Tranvías}
Hay 4 estaciones de tranvía repartidas por las esquinas del tablero. Si un 
jugador cae en una de estas casillas podrá moverse directamente a otra 
estación de tranvía de su elección pagando un precio por el desplazamiento.


\subsubsection{Párking}
Durante la partida, las multas, pagos de impuestos y demás ingresos 
a la banca en general,
se acumularán en un bote. El jugador que caiga en el parking se llevará
el bote completo. Si el bote está vacío, no ocurre nada.

\subsection{Preparación}
Cada jugador comienza con una cantidad de dinero determinada.
El turno se elige tirando un único dado, el que obtenga mayor valor empieza.
En caso de empate se vuelven a lanzar los dados para desempatar.
El orden de valores indica el orden de juego
de mayor a menor.


\subsection{Turno de juego}
Al comienzo del turno, el jugador lanzará los dados y moverá en función de lo 
que salga. Al caer en una casilla, se aplicará su efecto correspondiente.
Tras todos los efectos de casilla, el jugador podrá realizar acciones de 
gestión. Estas acciones son: comprar o vender construcciones, hipotecar o 
deshipotecar propiedades y comprar o vender propiedades a otros jugadores (intercambios, ver la Sección \ref{sec:intercambios}).

Haga lo que haga el jugador, su balance debe quedar positivo al final de su turno,
de lo contrario será declarado en bancarrota y eliminado de la partida.

\subsection{Subastas}
Las subastas serán pujas abiertas de valor oculto. Es decir, cada jugador
participante en la puja decidirá si participar o no, y en caso de participar
elegirá el valor por el que puja. El resto de jugadores no sabrá en ningún 
momento los jugadores que deciden participar en la puja ni los valores por los 
que pujan hasta el final de la subasta, donde se desvelarán los valores y 
el ganador de la puja. En caso de empate la propiedad no se dará a nadie.

Los jugadores que provocan una subasta no podrán participar en ella.

\subsection{Alquileres}
Si un jugador cae en una propiedad de otro jugador, deberá pagar el alquiler
correspondiente. El alquiler se calcula en función de si el dueño tiene el grupo 
de propiedades completo y de las construcciones que haya en la propiedad.

En caso de que el jugador no tenga dinero para pagar el alquiler al dueño,
el banco pagará el alquiler y al final del turno, el jugador deberá saldar 
su deuda con el banco mediante el hipotecado de propiedades o la venta de
construcciones o propiedades tanto al banco como a otros jugadores.

\subsection{Construcciones}
Una vez un jugador tiene el grupo completo de un color, podrá construir
casas en las propiedades de este grupo. El coste de construcción de cada 
casa vendrá marcado en cada propiedad. Las casas deben construirse de forma 
escalonada, es decir, si se quiere construir una segunda casa en una propiedad 
primero se tendrá que tener una casa en el resto de propiedades del grupo.

Una vez un jugador tenga 4 casas en todas las propiedades de un grupo, podrá
sustituir las 4 casas de una propiedad por un hotel por el precio de compra
de una casa.

A la hora de demoler construcciones se hará igualmente de forma escalonada en 
sentido inverso. En el caso particular de los hoteles, no se demolerán de golpe,
sino que primero se pasará de hotel a 4 casas.
El jugador recibirá la mitad del precio de compra de las construcciones demolidas.

\subsection{Intercambios}
\label{sec:intercambios}
Cada turno, el jugador podrá proponer intercambios a otros jugadores. Un 
intercambio se hará entre el jugador que lo propone y otro que podrá elegir.
Se podrán intercambiar dinero y cartas (propiedades, servidores y puentes)
por dinero y cartas. 

No se podrán intercambiar propiedades de un grupo de propiedades en el que
alguna de las propiedades tengan construcciones.

\subsection{Bonificaciones finales}
Al final de la partida, antes de desvelar al ganador, se aplicarán 3
bonificaciones finales aleatorias. Las bonificaciones darán una cantidad
de dinero determinada por un porcentaje del valor total de todos los jugadores 
sumado.
Estas bonificaciones vienen determinadas por estadísticas de la partida,
como el jugador que más se haya desplazado o el que más alquileres ha 
tenido que pagar al resto de jugadores, entre otras.

\subsection{Bancarrota}
Si un jugador no logra acabar su turno con balance positivo,
incluso después de hipotecar sus propiedades,
de vender todas sus construcciones y en general después de hacer 
gestiones para recaudar dinero,
todas las propiedades que le queden serán embargadas por el banco,
poniéndolas inmediatamente a la venta por su valor original.

Esto se aplicará también a los jugadores que se rindan, cuyo dinero en 
metálico quedará además absorbido por el banco sin ir al párking gratuito.



